%! Author = sriadityavedantam
%! Date = 3/25/26

\section{Open and Closed Sets}

\lesson{21}{Monday 29 September 2025 9:10}{Day 21}

\begin{definition}[$\eps$-Neighbourhood]
    Let $\eps > 0$.
    The neighborhood of $x\in\r$ is:
    \[
        \vee_{\eps}(x)\coloneqq \{y : \abs{x-y} < \eps\} = (x-\eps,x+\eps).
    \]
\end{definition}

\begin{definition}[Open Set]
    A set $U\subset\r$ is open if for every $x\in U$, there exists $\eps > 0$ such that $\vee_{\eps}(x)\subset U$.
\end{definition}

\begin{example}[Any open interval is open]
    If $x\in(a,b)$, take $\eps \coloneqq\min\{b-x, x-a\} > 0$.
    Then $\vee_{\eps}(x) \subset (a,b)$.
\end{example}

\begin{theorem}[Union of Open Sets]{
    The union of any family of open sets is open.
}
                 Let $\{U_\alpha\}_{\alpha\in I}$ be open and suppose $x\in\bigcup_{\alpha\in I} U_\alpha$.
                 Then $x\in U_{\alpha_0}$ for some $\alpha_0\in I$.
                 Since $U_{\alpha_0}$ is open, there exists $\eps > 0$ such that $\vee_{\eps}(x)\subset U_{\alpha_0}\subset\bigcup_{\alpha\in I} U_\alpha$.
\end{theorem}

\begin{example}[Examples of Open Sets]
    $\{0\}$ is not open.

    $(0,1]$ is not open.

    $\varnothing$ is open.
\end{example}

\begin{example}
    Let $I = \n$, and $\{q_1,q_2,\ldots\}$ be an enumeration of $\q$.
    For $n\in\n$, define
    \[
        U_n\coloneqq\vee_{1/2^n}(q_n).
    \]
    Then $U\coloneqq \bigcup_{n=1}^\infty U_n$ is open, and the total length satisfies
    \[
        \sum_{n=1}^\infty \frac{1}{2^{n-1}} = 4.
    \]
\end{example}

\begin{theorem}[Intersection of Open Sets]{
    The intersection of any finite collection of open sets is open.
}
                 Let $U_1,\ldots, U_N$ be open and $x\in\bigcap_{i=1}^N U_i$.
                 Then for each $i$, there exists $\eps_i > 0$ such that $\vee_{\eps_i}(x)\subset U_i$.
                 Let $\eps\coloneqq\min\{\eps_1,\ldots,\eps_N\} > 0$.
                 Then $\vee_{\eps}(x)\subset U_i$ for all $i$, hence $\vee_{\eps}(x)\subset\bigcap_{i=1}^N U_i$.
\end{theorem}

\begin{definition}[Limit Point]
    Let $E\subset\r$.
    A point $x_0\in\r$ is a limit point of $E$ if for every $\eps > 0$,
    \[
        (\vee_{\eps}(x_0)\cap E)\setminus\{x_0\} \neq \varnothing.
    \]
\end{definition}

\begin{example}
    If $E = (0,1)$, then every $x_0\in[0,1]$ is a limit point.
\end{example}

\begin{example}
    $\{c\}$ is a singleton set, which has no limit points.
\end{example}

\begin{definition}[Isolated Point]
    Given $E\subset\r$, $x_0\in E$ is an isolated point if it is not a limit point.
\end{definition}

\begin{proposition}{
    Given $E\subset\r$, $x_0\in\r$ is a limit point of $E$ if and only if there exists a sequence $(x_n)$ in $E\setminus\{x_0\}$ such that $x_n\to x_0$.
}
              Suppose $x_0$ is a limit point of $E$.
              For each $n\in\n$, choose $x_n\in \vee_{1/n}(x_0)\cap (E\setminus\{x_0\})$.
              Then $x_n\to x_0$.
              Conversely, if $x_n\to x_0$ with $x_n\in E\setminus\{x_0\}$, then every neighborhood of $x_0$ contains some $x_n$, hence $x_0$ is a limit point.
\end{proposition}

\begin{definition}[Closed]
    $E\subset\r$ is closed if it contains all its limit points.
\end{definition}

\begin{example}
    $(0,1)$ is not closed.

    $\{0\}$ is closed.

    $\varnothing$ is closed.

    $\z$ is closed.

    $\left\{\dfrac{1}{n} : n\in\n\right\}$ is not closed since $0$ is a limit point not contained in the set.
\end{example}

\begin{definition}[Closure]
    The closure of $E$, denoted $\overline{E}$ or $\cl(E)$, is the union of $E$ with all its limit points.
\end{definition}

\begin{example}
    $\overline{(0,1)} = [0,1]$.

    $\overline{\{0\}} = \{0\}$.
\end{example}

\begin{proposition}{
    $E$ is closed if and only if $E = \cl(E)$.
}
              If $E$ is closed, it contains all its limit points, so $\cl(E)=E$.

              Conversely, if $E=\cl(E)$, then $E$ contains all its limit points, hence is closed.
\end{proposition}